\section{Background and Related Work}

\paragraph{LLM Inference.}

LLM inference involves two stages: given a prompt $x_{\text{prompt}} = (x_1, x_2, \dots, x_m)$, the LLM first processes the prompt in parallel (prefill), then sequentially generates new tokens (decode), with each token $x_{t>m}$ conditioned on previously generated tokens:
$$
x_{t+1} = Sample(x | x_{1, \ldots, t}).
$$

In greedy sampling, the highest-probability token is selected iteratively until a stopping condition. Since each token depends on preceding outputs, generation is inherently sequential, requiring a separate forward pass per token, limiting throughput and underutilizing parallel hardware accelerators.

\paragraph{Speculative Decoding.}

Speculative decoding \citep{leviathan23} accelerates inference by generating multiple candidate tokens quickly using a lightweight model, which can then be verified in parallel by the primary LLM. The basic method has two core steps:

\begin{enumerate}
\item \emph{Speculation}: A smaller ``draft'' model rapidly produces speculative tokens $x_{\text{spec}} = (x_{t+1}, \dots, x_{t+n})$ based on the existing token prefix $x_{<t}$.

\item \emph{Verification}: The LLM verifies the draft tokens in parallel, accepting tokens up to the first discrepancy and discarding the rest.
\end{enumerate}

This shifts computation from sequential generation to parallel verification. However, draft models still require compute resources and add orchestration complexity. Model-based methods include Medusa~\citep{cai2024medusasimplellminference}, SpecInfer~\citep{specinfer} (which introduced tree-based speculation), multi-token prediction~\citep{gloeckle2024multitoken}, and blockwise parallel decoding~\citep{stern2018blockwise,kim2024accelerating}.

Recent \emph{model-free} methods include Prompt Lookup Decoding~\citep{saxena2023prompt}, Token Recycling~\citep{luo2024turningtrashtreasureaccelerating}, LLMA~\citep{yang2023inferencereferencelosslessacceleration}, and ANPD~\citep{ou2024losslessaccelerationlargelanguage}. These rely on small reference texts and lack adaptive speculation. SuffixDecoding is uniquely designed for agentic applications with long repetitive sequences.

\paragraph{Agentic AI Algorithms.}

Agentic applications structure tasks as multiple LLM calls, generating long repetitive token sequences. \emph{Self-consistency~\citep{wang2023selfconsistencyimproveschainthought}} samples multiple reasoning paths from the same prompt, sharing similar chain-of-thought steps. \emph{Self-refinement~\citep{madaan2023selfrefineiterativerefinementselffeedback}} iteratively fixes errors, revising small portions while preserving surrounding content. \emph{Multi-agent workflows~\citep{khot2023decomposedpromptingmodularapproach}} use specialized agents with narrow functions, producing repetitive outputs. These patterns create opportunities for exploiting long repeated sequences.

\paragraph{Methods for Accelerating LLM Agents.}

Recent works target latency in agentic applications. ALTO~\citep{10.1145/3642970.3655844} optimizes multi-agent workflows through pipelining and scheduling. Dynasor~\citep{fu2024efficientlyservingllmreasoning} early-terminates unlikely reasoning paths. SuffixDecoding takes an orthogonal speculative decoding approach and can be used in combination.

